\documentclass[11pt]{letter}
\usepackage[margin=1in]{geometry}
\usepackage{hyperref}

\signature{Brijesharun G \\ M.Tech (Software Engineering) \\ Dept. of Computing Technologies \\ SRM Institute of Science and Technology \\ Kattankulathur, India \\ brijesharun@gmail.com}

\date{\today}

\begin{document}

\begin{letter}{Editor-in-Chief}

\opening{Dear Editor,}

We are pleased to submit our manuscript entitled \textbf{``Automated Migration of Legacy Windows Desktop Applications to WinUI~3 Using Hybrid Rule-Based and Multi-Agent LLM Framework''} for consideration for publication in your journal.

This paper addresses a significant and previously unexplored problem in software modernization: the automated migration of legacy Windows desktop applications (WinForms, WPF, UWP) to Microsoft's modern WinUI~3 framework. While existing research has focused on programming language translation (e.g., C-to-Rust, PL/SQL-to-Java) or backend system modernization, no prior work addresses UI framework migration in the Windows desktop ecosystem.

Our key contributions include:

\begin{enumerate}
    \item A novel \textbf{hybrid migration framework} that combines deterministic rule-based transformations with a four-agent Large Language Model (LLM) architecture, where specialized agents handle analysis, translation, refactoring, and verification tasks collaboratively.

    \item A \textbf{Roslyn-based static analysis pipeline} for extracting UI control hierarchies, event bindings, and layout configurations from legacy source code into a framework-agnostic intermediate representation.

    \item A \textbf{curated dataset of 530 open-source Windows desktop applications} collected from public GitHub repositories, annotated with framework type, complexity tier, and file-level metadata.

    \item An \textbf{empirical evaluation} demonstrating that the hybrid multi-agent approach achieves 87.1\% compilation success rate, 97.0\% migration completeness, and 100\% UI parity---substantially outperforming rule-only (63.2\%) and single-agent LLM (76.2\%) baselines. The framework achieves a time reduction ratio of 5,867$\times$ over estimated manual effort.
\end{enumerate}

This work is particularly timely given the growing industry adoption of LLM-assisted software engineering workflows and Microsoft's ongoing push to modernize its desktop development platform. Our framework and dataset are publicly available at \url{https://github.com/floatingbrij/legacy-desktop-migration-llm} to support reproducibility.

This manuscript is original work, has not been published elsewhere, and is not under consideration by any other journal. All authors have read and approved the manuscript.

We suggest the following potential reviewers based on their expertise in related areas:

\begin{itemize}
    \item \textbf{Dr. Lakshmi Moti} --- Expert in multi-agent LLM architectures for code migration
    \item \textbf{Dr. Yuchen Luo} --- Researcher in LLM-assisted systems programming language migration
    \item \textbf{Dr. Jialun Cheng} --- Researcher in automated software modernization and refactoring
\end{itemize}

We look forward to your consideration.

\closing{Sincerely,}

\vspace{-0.3in}
\noindent \textbf{Corresponding Author:} \\
Dr. Hariprasad S \\
Dept. of Computing Technologies \\
SRM Institute of Science and Technology \\
Kattankulathur, India \\
haripras2@srmist.edu.in

\end{letter}
\end{document}
